\scanpage{18}
cette \unsure{dernière}, si $\operatorname{card}J=\operatorname{card}I-1=n-1$, on verra (par \unsure{récurrence}) que le cardinal de l'\unsure{ens} des classes est $(n-1)+1=n$.

Ceci posé, on voit d'abord que $E_i$ est dans l'une \inserted{(unique)} \unsure{de|des} ces $n$ classes, qui sont des ouverts. Soit $U_{J,i}$ la classe (ou ``composante'') qui contient $E_i$, et \unsure{considérons} les deux composantes $\Sigma_i^+,\Sigma_i^-$ de $X\setminus E_i$.

\noindent\begin{tabular}{@{}|p{.91\linewidth}@{}}
\uline{Lemme} Toute classe de $\Pi_{J\deleted{\unsure{i|j|,}}}^{\unsure{\circ|0}}$, distincte de $U_{J,i}$, est contenue dans $\Sigma_i^+$ ou $\Sigma_i^-$
\end{tabular}

C'est presque une tautologie : \unsure{Dire} \deleted{la} \inserted{qu'une} classe donnée \inserted{$U_{\unsure{\alpha|x}}$, de $\Pi_J^\circ$} ne contient pas $\deleted{\unsure{\Sigma}}\unsure{\overline{E_i}|E_i}$ i.e. est $\ne U_{J,i}$, signifie que il existe un $j\in J$ tel que le $\Sigma_j^+\in\pi_0(X\setminus \deleted{\unsure{\Sigma}}E_{\unsure{j|J}})$ \deleted{\unsure{choisi}} ``\unsure{composant|composante}'' de cette classe \deleted{est} est l'élément du $\pi_0$ qui ne contient pas $E_i$; il est donc contenu dans l'un des éléments \inserted{(soit $\Sigma_i^-$)} de $\pi_0(X\setminus\Sigma_i)$, d'après l'étude du cas \deleted{\unsure{deux}} \inserted{$\operatorname{card}I=2$}, \deleted{\unsure{OK|cela}} \unsure{â|à} fortiori $U_\alpha\subset\Sigma_i^-$, OK.
% archive sheet number: 17
